% 本文件由 LaTeX 排版 Agent 根据有效 translation.json 自动生成。
% author=Apocrypha; work_id=0832; title=关于圣母升天的伪经著作

\chapter{关于圣母升天的伪经著作}

% structure: p0001 source=h1 target=omitted reason=page-title-already-rendered-by-parent

% structure: p0002 source=h2 target=section reason=relative-heading-rank; base=h2

\section{圣史若望论天主之母安息记述}

当至圣荣耀的天主之母、永贞童贞玛利亚，按照她的习惯，前往我们主的神圣坟墓去焚香，并屈下圣膝，恳求那由她所生的基督——我们的天主——回到她那里时。犹太人看见她徘徊在神圣的坟墓旁，就来到司祭长跟前说：玛利亚每天去坟墓那里。司祭长召来他们派去守卫、不准任何人在圣墓前祈祷的卫兵，询问关于她的事，是否真是这样。卫兵回答说他们没有看见这样的事，因为天主没有允许他们看见她当时在那里。有一天，正是预备日，圣玛利亚照常来到坟墓；当她祈祷时，天开了，总领天使加俾额尔下来对她说：\Quote{万福！你生育了我们的天主基督！你的祈祷已上达诸天，蒙那由你而生的垂听；从今以后，照你的请求，你离开世界后，将到你儿子那里的天上之所，进入真正而永恒的生命。}\par

圣玛利亚听了圣总领天使的话后，回到圣白冷，与她同行的还有三位服侍她的童女。休息片刻后，她坐起来对童女们说：\Quote{拿香炉来，我要祈祷。}她们便遵命拿来了。她祈祷说：\Quote{我主耶稣基督，你因至高的美善屈尊由我而生，请垂听我的声音，派遣你的宗徒若望到我这里来，使我见到他便能分享喜乐；也请派遣你其余的宗徒，无论是已到你那里去的，还是仍在现世各处的，藉你的圣命，使他们无论在哪里都来；我见到他们，就要赞颂你极可称颂的圣名；因为我相信你在一切事上垂听你的婢女。}\par

当她祈祷时，我若望来了——圣神用云彩从厄弗所把我提去，放在我主母亲躺卧的地方。我进去到她身旁，赞颂那由她所生的，说：\Quote{万福，我主的母亲，你生育了我们的天主基督，欢喜吧，因为你正以极大的荣耀离开此生。}天主之母光荣天主，因为我若望到了她那里，记起主的话说：\BibleQuote{看，你的母亲}，以及，\BibleQuote{看，你的儿子}\BibleRef{若 19:26}。三位童女前来朝拜。天主之母对我说：\Quote{祈祷，并献香。}我便这样祈祷：\Quote{主耶稣基督，你行了奇妙的事，如今也在生你的母亲面前行奇妙的事；让你的母亲离开此生；并让那些钉死你而不信你的人蒙羞。}我结束祈祷后，圣玛利亚对我说：\Quote{把香炉拿给我。}她献了香，说：\Quote{光荣归于你，我的天主我的主，因为凡你在我升天之前所应许我的，如今在我身上都成全了，就是当我离开这世界时，你将与你的众天使带着荣耀降临到我这里。}我若望对她说：\Quote{我们的主耶稣基督，我们的天主，就要来了，你将要看见他，正如他应许你的。}天主之母回答我说：\Quote{犹太人已起誓，我死后他们要焚烧我的身体。}我回答她说：\Quote{你圣洁宝贵的身体决不会朽坏。}她又回答我说：\Quote{拿香炉来，献香，祈祷。}这时从天上传来声音，说\Quote{阿们。}我若望听见了这声音；圣神对我说：\Quote{若望，你听见了祈祷结束后在天空说话的声音吗？}我回答说：\Quote{是，我听见了。}圣神对我说：\Quote{你听见的这声音表示你的弟兄宗徒们和圣天使们就要来临，他们今天要来到这里。}\par

这时，我若望祈祷了。\par

圣神对宗徒们说：\Quote{你们众人一起，藉着云彩从世界各地而来，被旋风聚集在圣白冷，因我们的主耶稣基督的母亲之故；\Person{伯多禄}从\Place{罗马}，\Person{保禄}从\Place{提庇黎雅}，\Person{多默}从\Place{近印度}，\Person{雅各伯}从\Place{耶路撒冷}。\Person{安德肋}、伯多禄的兄弟，以及\Person{斐理伯}、\Person{路加}、\Person{西满·加纳人}，和已安眠的\Person{达陡}，被圣神从坟墓中唤醒；圣神对他们说：\Quote{不要以为现在是复活；但为此你们从坟墓中起来，好去问候我们的主救主耶稣基督的母亲，以尊崇和奇异之事，因为她离世升天之日临近了。}同样，\Person{马尔谷}也转道，从\Place{亚历山大}来到；他也与其余的人一起，如前所述，从各地而来。\Person{伯多禄}被云彩举起，站在天地之间，圣神使他稳固。同时，其余宗徒也被云彩提去，与\Person{伯多禄}一同出现。因此，如上所述，他们都在圣神带领下聚集了。}\par

我们进去到我们主和天主的母亲身旁，朝拜后，说：\Quote{不要害怕，也不要悲伤；那由你而生的天主上主，将带着荣耀把你带出这个世界。}她因天主她的救主而欢喜，在床上坐起来，对宗徒们说：\Quote{现在我已相信我们的主和天主将从天降来，我将看见他，并这样离开此生，如同我看见你们来了一样。我希望你们告诉我，你们怎么知道我要离去而来到我这里，以及你们从哪些国家、经过多远来到这里，竟如此急忙前来探望我。因为那由我而生的我们的主耶稣基督，万有的天主，并没有隐藏这事；因为我现在仍深信他是至高者的儿子。}\par

\Person{伯多禄}回答，对宗徒们说：\Quote{我们各人，按照圣神所宣告和命令我们的，将详细情况告知我们主的母亲。}\newline{} \Person{若望}回答说：\Quote{正当我进入\Place{厄弗所}的圣祭台举行神圣礼仪时，圣神向我说：\InnerQuote{你们的主的母亲离世的时候近了；你去\Place{白冷}向她请安。}有一朵光云攫取了我，将我放在你们躺卧的门前。}\newline{} \Person{伯多禄}也回答说：\Quote{我住在\Place{罗马}，将近黎明时，通过圣神听到一个声音向我说：\InnerQuote{你们的主的母亲要离世了，因为时候近了；你去\Place{白冷}向她请安。}看，一朵光云攫取了我；我也看见其他宗徒乘着云彩向我而来，有一个声音向我说：\InnerQuote{你们都去\Place{白冷}。}}\newline{} \Person{保禄}也回答说：\Quote{我住在离\Place{罗马}不远的城市，称为\Place{提庇黎雅}地区，听到圣神向我说：\InnerQuote{你们的主的母亲已离开此世，正通过离世向天上行进；但你也去\Place{白冷}向她请安。}看，一朵光云攫取了我，将我放在与你们相同的地方。}\newline{} \Person{多默}也回答说：\Quote{我穿越\Place{印度}地区时，因基督的恩宠，宣讲大行其道，国王的姐妹的儿子\Person{拉布达努斯}，即将在宫殿里由我\OriginalTerm{施行印记}{sealed}，突然圣神向我说：\InnerQuote{多默，你也去\Place{白冷}向你们的主的母亲请安，因为她正离世往天上去了。}一朵光云攫取了我，将我放在你们旁边。}\newline{} \Person{马尔谷}也回答说：\Quote{当我在\Place{亚历山大}城完成\Emph{第三天}的礼仪正典时，正当我祈祷，圣神攫取了我，将我带到你们这里。}\newline{} \Person{雅各伯}也回答说：\Quote{当我在\Place{耶路撒冷}时，圣神命令我说：\InnerQuote{去\Place{白冷}，因为你们的主的母亲正离世。}看，一朵光云攫取了我，将我放在你们旁边。}\newline{} \Person{玛窦}也回答说：\Quote{我曾光荣，并继续光荣天主，因为当我在船上被风暴袭击时，海面波涛汹涌，突然一朵光云笼罩了风暴的巨浪，使它变为平静，并攫取了我，将我放在你们旁边。}\newline{} 先前到达的人也照样回答，说明了他们如何到来。\Person{巴尔托禄茂}说：\Quote{我在\Place{特拜斯}宣讲圣言时，看，圣神向我说：\InnerQuote{你们的主的母亲正离世；所以你去\Place{白冷}向她请安。}看，一朵光云攫取了我，将我带到你们这里。}\par

宗徒们将这一切事，即他们为何而来以及如何而来，都告诉了圣洁的天主之母；她举手向天祈祷说：\Quote{主，我朝拜、赞美并光荣祢极可赞美的名，因为祢眷顾了祢婢女的卑微，因为全能者为我行了大事；看，从今以后万代都要称我有福。}\BibleRef{路 1:48}祈祷后，她对宗徒们说：\Quote{献香并祈祷。}他们祈祷后，从天上传来雷声，并有威严的声音，仿佛战车之声；看，有一大队天使和异能者，并似人子的声音被听见，色辣芬环绕着圣洁无玷的天主之母和童贞所躺卧的房子，以致所有在\Place{白冷}的人都看见了这一切奇事，并来到\Place{耶路撒冷}报告所发生的一切奇事。当这声音被听见时，太阳和月亮突然出现在房子周围；一群首生的圣人站在主的母亲躺卧的房子旁边，为她的荣耀和尊贵。我也看见许多征兆发生：盲者看见，聋者听见，瘸子行走，癞病人洁净，被不洁之神附身的人得到痊愈；每一个患病或抱恙的人，只要触摸她躺卧的房子的外墙，就呼喊说：\Quote{圣玛利亚，生了基督我们天主的主，求祢垂怜我们。}他们立刻被治愈。有许多来自各地为了祈祷而住在\Place{耶路撒冷}的人，听说了通过主的母亲在\Place{白冷}所行的征兆，就来到这地方寻求各种疾病的治愈，也得到了。那一天，在痊愈的人群中和旁观者中，有说不出的喜乐，他们都光荣基督我们的天主和祂的母亲。整个\Place{耶路撒冷}连同\Place{白冷}都因着圣咏和灵歌而庆祝。\par

犹太人的司祭同他们的百姓对所发生的事感到惊奇；但因极度仇恨的驱使，又因轻浮的推论，他们集会决定派遣人去攻击在那里（\Place{白冷}）的天主之母和圣宗徒。于是，犹太民众向\Place{白冷}进发，当行至一英里处时，他们看见了一个可怕的异象，他们的脚被钉住；此后他们回到同胞那里，将一切可怕的异象报告给司祭长。他们更加怒火中烧，去见总督，喊叫说：\Quote{犹太民族被这女人毁灭了；把她从\Place{白冷}和\Place{耶路撒冷}省赶走。}总督对这一切奇事感到惊讶，对他们说：\Quote{我既不会把她从\Place{白冷}赶走，也不会从任何其他地方赶走。}犹太人继续喊叫，并指着\Person{提庇留·凯撒}的健康起誓，要求把宗徒们从\Place{白冷}带出来。\Quote{你若不然，我们就向凯撒报告。}于是，他被逼无奈，派一名千夫长率兵去\Place{白冷}攻打宗徒。圣神对宗徒和主的母亲说：\Quote{看，总督派了一名千夫长来攻打你们，犹太人已经骚动。因此，离开\Place{白冷}吧，不要害怕；因为看，我要用云彩将你们带到\Place{耶路撒冷}；因为圣父、圣子、圣神的权能与你们同在。}宗徒们遂立即起身，走出屋子，抬着夫人、天主之母的床，向\Place{耶路撒冷}发奋。立刻，正如圣神所说，他们被云彩举起，就发现已在\Place{耶路撒冷}夫人的家中。他们起来，连续五天不停地赞美歌唱。千夫长来到\Place{白冷}，在那里既找不到主的母亲，也找不到宗徒，就抓住\Place{白冷}人，对他们说：\Quote{你们不是去告诉总督和司祭一切发生的征兆和奇事，以及宗徒们从各地出来的事吗？那么，他们在哪里？来，到\Place{耶路撒冷}的总督那里去。}因为千夫长不知道宗徒和主的母亲已经去了\Place{耶路撒冷}。千夫长于是带着\Place{白冷}人进去见总督，说他一个人也没找到。五天后，总督、司祭和全城的人都从那里发生的征兆和奇事得知，主的母亲和宗徒们在\Place{耶路撒冷}自己的家中。许多男人、妇女和贞女聚集起来，喊叫说：\Quote{圣童贞，你生了我们的天主基督，不要忘记人类。}这些事发生时，犹太民众连同司祭们更加被仇恨驱动，拿了木柴和火，前来要烧毁主的母亲与宗徒们居住的房子。总督站在远处观看这情景。犹太民众来到屋门前时，看，忽然从里面藉着天使发出一股火的能力，烧死了许多犹太人。全城甚为恐惧；他们光荣了由她而生的天主。总督看见所发生的事，就向全体民众喊道：\Quote{真正地，那由童贞所生、你们想赶走的那一位，是天主之子；因为这些征兆是属真天主的。}犹太人中间起了分裂；许多人因所发生的征兆而相信了我们主耶稣基督的名。\par

天主之母、终身童贞玛利亚，主的母亲，经历了这一切奇妙之事后，我们宗徒与她同在\Place{耶路撒冷}时，圣神对我们说：你们知道，在主日，总领天使\Person{加俾额尔}向童贞玛利亚传报喜讯；在主日，救主诞生于\Place{白冷}；在主日，\Place{耶路撒冷}的儿童手持棕榈枝出来迎接祂，说：\Quote{贺三纳于至高之天，奉主名而来的，当受赞颂。}；在主日，祂从死者中复活；在主日，祂将审判生者死者；在主日，祂将从天上降临，为荣耀和尊荣那生育祂的圣洁荣耀童贞女之离世。在同一主日，主的母亲对宗徒说：\Quote{献上乳香，因为基督正偕同天使大军来临；看，基督已近在咫尺，坐在\OriginalTerm{革鲁宾}{cherubim}的宝座上。}当我们众人祈祷时，出现了无数天使，主以大能乘驾\OriginalTerm{革鲁宾}{cherubim}；看，一道光芒射向圣童贞，因她的独生子临在，天上万有都俯伏朝拜祂。主对祂的母亲说：\Quote{玛利亚。}她回答说：\Quote{主，我在这里。}主对她说：\Quote{不要悲伤，让你的心欢喜快乐；因为你已蒙恩宠，得以看见我父赐给我的荣耀。}天主圣母亲眼观看，见祂身上有一种人言无法述说或理解的荣耀。主停留在她身旁，说：\Quote{看，从今以后，你宝贵的身体将被移至乐园，你的圣灵魂将被带到天上我父的宝库中，在极其光明之处，那里有平安与圣天使的喜乐——还有其他。}主的母亲回答祂说：\Quote{主，请将你的右手放在我身上，祝福我。}主伸出祂无玷的右手，祝福了她。她握住祂无玷的右手，亲吻它，说：\Quote{我朝拜这只右手，它创造了天地；我呼求你当受赞美之名，基督，天主，万世之王，圣父的独生子，收纳你的婢女，你曾屈尊由我卑微之所生育，以拯救人类，藉你那不可言喻的救世计划；求你赐助凡呼求、祈祷或提说你婢女之名的每一个人。}她这样说时，宗徒们走到她脚前，俯伏朝拜说：\Quote{主的母亲，请给世界留下祝福，因为你即将离开它。你已祝福了世界，并在它沉沦时高举了它，因你生了世界之光。}主的母亲祈祷，在祈祷中这样说：\Quote{天主，你因你的大慈大悲，从天上派遣你的独生子居住在我卑微的身体内，你曾屈尊由我\Emph{这样}卑微之人所生，求你怜悯世界，以及凡呼求你名的每个灵魂。}她又祈祷说：\Quote{主，天上的君王，永生天主子，求你收纳凡呼求你名的每一个人，使你的诞生得荣耀。}她又祈祷说：\Quote{主耶稣基督，你在天上地下拥有全部权能，在此恳求中，我呼求你圣名；在任何时间、任何地方，凡提及我名之处，求你圣化那地方，并光荣那些因我名而光荣你的人，收纳他们的一切奉献、一切恳求和一切祈祷。}她这样祈祷后，主对祂的母亲说：\Quote{让你的心欢喜快乐；因为每一种恩惠和每一种恩赐都已从我天上的父那里、从我、并从圣神赐给了你：凡呼求你名的灵魂，必不蒙羞，反要在现今的世界和未来的世界，在我天父面前，获得怜悯、安慰、支持和信心。}主转身对\Person{伯多禄}说：\Quote{时辰已到，开始唱赞美诗吧。}\Person{伯多禄}便开始唱赞美诗，天上万有都回应阿肋路亚。那时，主的母亲面容比光明更耀眼，她起身亲手祝福了每一位宗徒，众人都光荣了天主；主伸出祂无玷的双手，接了她的圣洁无玷的灵魂。在她无玷灵魂离世时，那地方充满了馨香和不可言喻的光明；看，有声音从天上传来：\BibleQuote{你在妇女中应受赞颂。}\Person{伯多禄}、我\Person{若望}、\Person{保禄}和\Person{多默}跑去为她宝贵的脚裹上圣布；十二位宗徒将她宝贵的圣体放在床上，抬着前行。看，他们正抬着时，有一个出身高贵的希伯来人，名叫\Person{耶福尼雅}，冲撞遗体，将手放在床上；看，主的一位天使以无形之力，用火剑将他的双手从肩膀上砍下，使它们悬在床上，悬在空中。看到所发生的这个奇迹，所有在场的犹太人都喊道：\Quote{真正由你生育的，是天主，哦，天主之母、终身童贞玛利亚。}至于\Person{耶福尼雅}本人，当\Person{伯多禄}命令他，为彰显天主的奇事时，他站起来，在床后喊道：\Quote{圣玛利亚，你生了基督，祂是天主，求你怜悯我。}\Person{伯多禄}转身对他说：\Quote{借着由她所生者的名，你那被砍去的双手必将重新接上。}立刻，因\Person{伯多禄}的话，悬在圣母床上的双手飞下来，接在\Person{耶福尼雅}身上。他便相信了，并光荣了由她所生的基督天主。\par

当这奇迹完成后，宗徒们抬起床榻，将圣母珍贵圣洁的遗体安放在\Place{革责玛尼}的一个新坟墓中。看，从我们的天主之母的圣墓中发出芬芳的香气；三日之久，听见无形天使的歌声，荣耀那由她所生的基督我们的天主。第三天结束时，歌声不再听见；从那时起，众人都知道她那无玷珍贵的身体已被迁移到天堂。\par

在它被迁移之后，看，我们看见\Person{圣若翰洗者之母依撒伯尔}，以及\Person{圣母之母亚纳}，还有\Person{亚巴郎}、\Person{依撒格}、\Person{雅各伯}、\Person{达味}，唱着阿肋路亚，以及诸圣的唱诗班朝拜主之母的圣髑，那地方充满光明，再没有什么光比这光更辉煌，并且在她珍贵圣洁的身体被迁移到天堂的那地方，有丰盛的香气，以及赞美由她所生者的旋律——甜美的旋律，令人永不满足，这种旋律只赐予贞女，唯独她们才能听见。因此，我们宗徒目睹了她圣洁身体的突然珍贵的迁移，光荣了天主，祂在我们主耶稣基督之母离世时向我们显示了祂的奇事；愿我们众人因她的庇护、支持和保护，在现世和来世，都堪当领受她的祈祷和转求，时时处处光荣她的独生子，与圣父及圣神一起，直到永远。阿们。\par

% structure: p0015 source=h2 target=section reason=relative-heading-rank; base=h2

\section{圣母安息}

% structure: p0016 source=h3 target=subsection reason=relative-heading-rank; base=h2

\subsection{拉丁文第一式}

\Emph{关于童贞玛利亚的安息}\par

在主受难之前的那段时间，在圣母向圣子提出的许多问题中，她开始询问关于自己的离世，对他这样说：——\Quote{最亲爱的儿子，我恳求祢的圣洁，当我的灵魂离开身体时，祢让我在三天前知道；并求祢，爱子，与祢的天使一同接纳它。}于是祂接受了祂爱母的祈祷，对她说：\Quote{永生天主的宫殿和圣殿，有福的母亲，诸圣之后，在众女人中受赞颂者，在祢将我还怀在祢的胎中之前，我一直守护着祢，并每日以我的天使食物喂养祢，如祢所知道的，在祢怀了我、哺养了我、带我逃往埃及、为我忍受了许多艰辛之后，我怎能离弃祢呢？因此，要知道，我的天使们一直守护着祢，并将守护祢直到祢离世。但在我为人类受难之后，正如经上所记载的，第三天复活，四十天后升天，那时祢将看见我与众天使、总领天使、圣人、贞女以及我的门徒们一同来临时，祢便可确知祢的灵魂将离开身体，我将把它带到天堂，在那里它永不会遭受任何磨难或痛苦。}于是她欢欣荣耀，亲吻了她儿子的膝盖，并赞美了天地万物的创造者，因祂藉她儿子耶稣基督赐给她如此恩赐。\par

在第二年后，就是在我们的主耶稣基督升天之后，最蒙福的童贞玛利亚日夜持续祈祷。在她安息前的第三天，主的天使来到她面前，问候她说：\Quote{万福，玛利亚，充满恩宠！主与你同在。}她回答说：\Quote{感谢天主。}他又对她说：\Quote{接过这枝主应许给祢的棕榈枝。}她感谢天主，极其喜乐地从天使手中接过了给她的棕榈枝。主的天使对她说：\Quote{祢的安息将在三天后。}她回答说：\Quote{感谢天主。}\par

于是她叫来了\Person{阿黎玛特雅}人\Person{若瑟}，以及主的其他门徒。当他们都聚集在一起时，在场的亲戚和熟人，她向众人宣布了她的离世。然后蒙福的玛利亚沐浴，如王后般装扮，等待她儿子的来临，正如祂所应许的。她请求所有亲戚留在她身边，安慰她。她身边还有三个贞女：\Person{塞福辣}、\Person{阿彼盖耳}和\Person{匝厄耳}。但我们主耶稣基督的门徒们早已分散到世界各地向天主子民传道。\par

到了第三时辰，有大雷、大雨、闪电、灾难和地震，此时圣母玛利亚正站在她的内室里。圣史兼宗徒\Person{若望}忽然从\Place{厄弗所}被带来，进入了圣母玛利亚的内室，向她请安，对她说：\Quote{万福，玛利亚，满被恩宠！主与你同在。}她回答说：\Quote{感谢天主。}她起身，亲吻了圣\Person{若望}。圣母玛利亚对他说：\Quote{我最亲爱的儿子，你为何在这样的时刻离开我，没有留意你师傅的命令，就是祂悬在十字架上时吩咐你照顾我的？}他屈膝求饶。于是圣母玛利亚祝福了他，又亲吻了他。当她正要问他从何处来，以及为何来\Place{耶路撒冷}时，看，所有主的门徒，除了称为狄狄摩的\Person{多默}，都被一朵云彩带到了圣母玛利亚内室的门口。他们站住，走了进去，用以下的话向圣母请安，并敬礼她：\Quote{万福，玛利亚，满被恩宠！主与你同在。}她急忙起身，俯身，亲吻他们，并感谢天主。以下是那些被云彩带到那里的主的门徒的名字：圣史\Person{若望}和他的兄弟\Person{雅各伯}，\Person{伯多禄}和\Person{保禄}，\Person{安德肋}，\Person{斐理伯}，\Person{路加}，\Person{巴尔纳伯}，\Person{巴尔多禄茂}和\Person{玛窦}，称为义者的\Person{玛弟亚}，加纳人\Person{西满}，\Person{犹达}和他的兄弟，\Person{尼苛德摩}和\Person{马克西米安}，以及其他许多无法计数的人。于是圣母玛利亚对她的弟兄们说：\Quote{这是什么事？你们全都来到了\Place{耶路撒冷}？}\Person{伯多禄}回答说：\Quote{我们需要问你这件事，你反倒问我们？确实，我想，我们谁也不知道为何今天如此迅速地来到这里。我刚才还在\Place{安提约基雅}，现在却在这儿了。}大家都明白地说出他们当天所在的地方。他们听了这些事后，都惊异自己来到了这里。圣母玛利亚对他们说：\Quote{我曾在祂受难之前求过我的圣子，要祂与你们都在我临终时临在；祂赐给了我这份恩赐。因此你们可以知道，我的离世将在明天。你们要与我一起醒寤祈祷，这样当主来收取我的灵魂时，祂会发现你们醒寤。}于是众人都许诺要醒寤。他们整夜醒寤祈祷，咏唱圣咏和圣歌，大放光明。\par

到了主日，第三时辰，正如圣神曾以云彩降临在宗徒们身上一样，基督也带着众多天使降来，接收了祂爱母的灵魂。因为那时有如此的光辉和甜蜜的芬芳，天使们咏唱着雅歌（主曾说：\Quote{如百合花在荆棘中，我的爱人在女子中亦然}，\BibleRef{歌 2:2}），所有在场的人都俯伏于地，就如基督在\Place{大博尔山}上在他们面前显圣容时宗徒们俯伏一样，整整一个半小时无人能起来。但当那光离去时，与光同时，蒙福童贞玛利亚的灵魂被圣咏、赞美诗和雅歌带到天上。云彩上升时，全地震动，瞬间\Place{耶路撒冷}的所有居民都公开目睹了圣玛利亚的离世。\par

就在那同一时刻，撒殚进入了他们心中，他们开始考虑如何处置她的身体。他们拿起武器，要焚烧她的身体并杀害宗徒，因为从她那里产生了以色列因罪恶而流离失所、以及外邦人的聚集。但他们被击打而盲目，头撞墙壁，互相打击。于是宗徒们被如此的光明所惊，起身，咏唱着圣咏，将圣身从\Place{熙雍山}抬下到\Place{约沙法特山谷}。当他们走到路中间时，看，有一个名叫\Person{勒乌本}的犹太人，想要把载有圣母玛利亚身体的抬架扔到地上。但他的双手枯干，直到手肘；他无论愿意与否，一直下到\Place{约沙法特山谷}，哭泣哀叹，因为他的手被举向抬架，无法收回。他开始恳求宗徒们，藉他们的祈祷使他得救，成为基督徒。于是宗徒们屈膝，求主释放他。他立即在那时刻痊愈了，感谢天主，并亲吻了诸圣及宗徒之后的脚，就在那地方受了洗，开始宣扬我们的天主耶稣基督的名。\par

于是宗徒们极其恭敬地将身体安放在坟墓中，因着极大的爱情和甜蜜而哭泣歌唱。忽然，有光从天上照耀他们周围，他们仆倒在地，圣身被天使举升到天上。\par

那时，最蒙福的多默突然被提到橄榄山上，看见那最蒙福的身体升天，便大声呼喊说：\Quote{哦，圣善的母亲，蒙福的母亲，无玷的母亲，若我如今因看见你而蒙恩，请以你的怜悯使你的仆人喜乐，因为你正升天。}接着，宗徒们曾环绕最圣洁身体的那条腰带从天上被掷下，落在蒙福的多默身上。他拿起它，亲吻它，感谢天主，又回到约沙法特谷。他找到所有宗徒和一大群人，他们因所见的光明而捶胸。蒙福的伯多禄看见他，彼此亲吻后，对他说：\Quote{你确实总是固执而不信，因为你的不信，天主不喜悦你与我们一同为救世主之母安葬。}他捶着胸说：\Quote{我知道并坚信，我一直是个邪恶不信的人；因此，我为我的固执和不信仰请求你们所有人的宽恕。}他们便都为他祈祷。然后蒙福的多默说：\Quote{你们将她的身体放在了哪里？}他们用手指指着坟墓说：\Quote{那被称为最圣洁的身体不在那里。}蒙福的伯多禄对他说：\Quote{你曾有一次不信我们的主和老师复活的话，除非你亲手触摸他、看见他；如今你怎能相信我们，说圣洁的身体在这里？}他仍坚持说：\Quote{不在这里。}他们便仿佛恼怒地走到坟墓那里——那是一个在岩石中凿成的新坟墓——搬开了石头，却没有找到身体，不知该说什么，因为他们被多默的话证实了。于是蒙福的多默告诉他们，他如何在印度举行弥撒——他仍穿着司铎礼服。他因不知天主的圣言，而被提到橄榄山上，看见蒙福的玛利亚的最圣洁身体升天，便祈求她赐福。她听了他的祈祷，将她的腰带扔给了他。宗徒们看见那条曾系在她身上的腰带，光荣天主，都因蒙福的玛利亚赐予多默的祝福以及他看见那最圣洁身体升天而请求多默宽恕。蒙福的多默便祝福了他们，说：\BibleQuote{看，弟兄们同居共处，是多么美好，多么愉快！}\par

那曾带他们来的云彩又把他们每个人带回自己的地方，正如斐理伯为太监施洗时所记载的，如宗徒大事录所读的\BibleRef{宗 8:39}；又如哈巴谷先知将食物带给在狮子坑中的达尼尔，然后迅速回到犹太。同样，宗徒们迅速回到他们最初所在之处，向天主的人民宣讲。这并不希奇，因为那进入童贞女又在她胎闭的状态中出来、在门关闭时进入门徒中间\BibleRef{若 20:19}、使聋子听见、复活死人、洁净癞病人、使瞎子看见、并行了其他许多奇事的天主，能行这样的事。相信这事并无可疑之处。\par

我是若瑟，曾将主的身体安放在我的坟墓中，并看见他复活；在主升天之前和之后，我一直守护着他最神圣的殿堂——蒙福的终身童贞玛利亚，并将那从天主口中说出的事和以上所记按天主的审判所成就的事，以文字并牢记于心。我向众人——犹太人和外邦人——明示了我亲眼所见、亲耳所闻的一切；只要我一息尚存，我必不止息地宣告它们。让我们不断地恳求她——她的升天今日普世尊崇敬礼——在天上在她最仁慈的圣子面前记念我们；愿光荣与赞颂归于他，至于无穷世之世。阿们。\par

% structure: p0028 source=h3 target=subsection reason=relative-heading-rank; base=h2

\subsection{第二种拉丁文形式}

在此开始蒙福童贞玛利亚的安眠。\par

1. 因此，当主和救世主耶稣基督为整个世界获得生命而被钉在十字架上悬于木架时，他看见他的母亲站在十字架旁，以及他所特别爱的那位宗徒若望——因为他与众宗徒不同，是身体贞洁的童贞者。于是他托付了圣玛利亚给若望，对他说：\BibleQuote{看，你的母亲！}又对她说：\BibleQuote{看，你的儿子！}\BibleRef{若 19:26}从那时起，天主之母在世上居住的日子里，特别受若望照管。宗徒们拈阄分派天下传教后，她便定居在若望的父母家里，靠近橄榄山。\par

2. 因此，在\Person{基督}击败死亡并升到天上后的第二年，有一天，\Person{玛利亚}因渴望\Person{基督}而心中炽热，开始独自在她住所的遮蔽下哭泣。看，一位身穿极其光明之衣的天使站在她面前，说出问候的话说：\Quote{万福！上主所祝福的，请接受那藉其先知命令雅各伯得平安者的问候。看，祂说，一根棕榈枝——我从\OriginalTerm{上主的乐园}{paradise of the Lord}带来给你——你将在第三日被接离身体时，让人在你灵柩前抬着它。因为，看，你的儿子正以宝座、天使和天上万军等候你。}于是\Person{玛利亚}对天使说：\Quote{我恳求\Person{主耶稣基督}的所有宗徒都聚集到我这里来。}天使对她说：\Quote{看，今日因我\Person{主耶稣基督}的能力，所有宗徒都会来到你这里。}\Person{玛利亚}对他说：\Quote{我求你赐福于我，使我在灵魂离开身体的那时刻，没有\OriginalTerm{下界的势力}{power of the lower world}能抵挡我，并且使我不看见\OriginalTerm{黑暗之君}{prince of darkness}。}天使说：\Quote{\OriginalTerm{下界的势力}{power of the lower world}确实不会伤害你；你的主天主，我是祂的仆人和使者，已赐给你永恒的祝福；但不要以为不看见\OriginalTerm{黑暗之君}{prince of darkness}的特权是由我赐给你的，而是由你怀中的那一位；因为万有永远属于祂。}说了这话，天使带着极大的光辉离去。那棕榈枝闪耀着极其明亮的光。于是\Person{玛利亚}脱下衣服，穿上更好的衣裳。她拿着从天使手中接过的棕榈枝，出去到\Place{橄榄山}上，开始祈祷说：\Quote{主啊，除非你怜悯了我，我不配怀你；但我仍保存了你托付给我的珍宝。因此，\OriginalTerm{荣耀的君王}{King of glory}，我求你，不要让\OriginalTerm{革赫纳}{Gehenna}的势力伤害我。因为如果诸天和天使每日在你面前战栗，那从土所造的人岂不更当如此？人本身一无所有，只从你仁慈的恩惠中领受。主啊，你是永远当受赞美的天主。}说了这话，她返回自己的住所。\par

3. 看，忽然，当\Person{圣若望}在\Place{厄弗所}传道时，在主日白天第三时辰，发生了大地震，一朵云彩将他举起，从众人眼前带离，送到\Person{玛利亚}所住房屋的门前。他敲门后立即进去。\Person{玛利亚}看见他，喜乐洋溢，说：\Quote{我求您，我儿若望，要记住我\Person{主耶稣基督}托付你照顾我的话语。因为，看，第三天我将离世，我听到了犹太人的计划，他们说：\InnerQuote{我们等待那生育那诱惑者的妇人死的那天，然后用火烧她的身体。}}于是她叫了\Person{圣若望}，领他进入房子里的内室，将她的葬衣和那从天使手中接过的光耀棕榈枝给他看，吩咐他应在她的床榻被送往坟墓时让人抬在床前。\par

4. \Person{圣若望}对她说：\Quote{我怎能独自为你举行葬礼，除非我的弟兄和\Person{主耶稣基督}的同工宗徒们来向你的遗体致敬？}看，忽然，因天主的命令，所有宗徒都被提起，被云彩托起，从他们宣讲天主圣言的地方，被放在\Person{玛利亚}所住房屋的门前。他们互相致意，感到惊奇，说：\Quote{主为何把我们聚集在这里？}\par

5. 于是所有宗徒同心合意地喜乐，完成了祈祷。他们说了\Quote{阿们}后，看，忽然有\Person{真福若望}前来，将这一切事告诉了他们。众宗徒于是进入屋子，找到了\Person{玛利亚}，向她致意说：\Quote{你因创造了天地的上主而蒙受祝福。}她对他们说：\Quote{愿你们平安，最亲爱的弟兄们！你们是怎样来到这里的？}他们向她述说他们怎样各自被天主的圣神借云彩提起，降在同一个地方。她对他们说：\Quote{天主没有使我见不到你们。看，我将走世人必经之路，我不怀疑主如今引领你们到此，是为了在我即将来临的痛苦中给我安慰。因此现在我恳求你们，要毫不间断地同心合意警醒，直等到主来的那时刻，那时我将离世。}\par

6. 他们围坐一圈安慰她，当他们在赞美天主中度过了三天后，看，第三天，大约白天第三时辰，一种沉睡抓住了那屋里所有的人，除了宗徒们和那里的三位童贞女，没有人能保持清醒。看，忽然\Person{主耶稣基督}与一大群天使一同来临；极大的光明降临在那地方，天使们唱着赞美诗，共同赞美天主。于是\Person{救主}说道：\Quote{来，最宝贵的珍珠，进入\OriginalTerm{永生之所}{receptacle of life eternal}。}\par

7. 于是玛利亚俯伏在地，朝拜天主，说：\Quote{愿祢光荣的圣名受赞美，我主、我的天主，祢竟屈尊拣选祢的婢女，并将祢隐藏的奥秘托付给我。因此，光荣的君王，求祢记念我，因为祢知道我用全心爱了祢，并守护了托付给我的珍宝。求祢收纳祢的仆人，救我脱离黑暗的权势，免得撒但的攻击阻挡我，也叫我不致看见污秽的灵挡在我的路上。}救主回答她说：\Quote{当我受圣父派遣，为拯救世界而被悬在十字架上时，黑暗之君来到我面前；但他在我身上找不到他作为的任何痕迹\BibleRef{若 14:30}，便败退被践踏。但当你看见他时，你因人类的律法而确实会看见他，因为你已到了生命的终点；但他不能伤害你，因为我与你同在，要帮助你。安心去吧，因为天军正等候你，要领你进入乐园的喜乐。}主这样说后，玛利亚从地上起来，斜靠在床上，感谢天主，便交付了灵魂。宗徒们看见她的灵魂如此洁白，没有凡人所能言说；因为它极度的光芒超越了雪、一切金属和闪亮的白银的洁白。\par

8. 于是救主说：\Quote{伯多禄，起来，拿起玛利亚的身体，送到城东的右边去，你会在那里找到一座新墓穴，把她安放在里面，等候我到你们那里去。}说完这话，主便将圣玛利亚的灵魂交付给弥额尔——他是乐园的统治者和犹太民族的首领；加俾额尔与他们同行。救主随即与天使一同被接升天。\par

9. 当时在场的三位童女抬起真福玛利亚的身体，要按殡葬礼仪为她洗濯。她们脱去她的衣服后，那神圣的身体发出如此光辉，以致虽可触摸以准备埋葬，但因那极强的闪光，它的形状却无法看见；只是主的光辉显为大，什么也看不见，那身体洗濯后完全洁净，未染任何污秽湿气。给祂穿上寿衣后，那光渐渐暗了下来。真福玛利亚的身体如同百合花，且发出极大的芳香，没有别的芳香可比拟。\par

10. 于是宗徒们把圣尸放在灵床上，彼此说：\Quote{谁要在她的灵床前持这棕榈枝？}若望对伯多禄说：\Quote{你在宗徒职务中居我们之先，应在她的床前持这棕榈枝。}伯多禄回答说：\Quote{你是主从我们中所选的唯一童身，你得了极大的恩宠，曾靠在祂的胸膛\BibleRef{若 13:23}。祂为我们的救恩悬在十字架上时，亲口把她托付给你。因此你应当持这棕榈枝，让我们抬起那身体，抬到埋葬之地。}此后，伯多禄举起身体说：\Quote{抬起身体}，开始歌唱说：\Quote{以色列出了埃及。阿肋路亚。}其他宗徒与他一同抬起真福玛利亚的身体，若望持着发光的棕榈枝走在前头。其他宗徒以极甜美的声音歌唱。\par

11. 看，新的奇迹。灵床上方出现一片极大的云彩，如同常出现在月亮光辉旁的大圆环；云中有天使军队发出甜美的歌声，大地因那极甜美之声而回响。于是约有一万五千人从城中出来，惊奇说：\Quote{这极大的甜美声音是什么？}有一个人站起来对他们说：\Quote{玛利亚已离世，耶稣的门徒正在她周围唱赞美诗。}他们观看，见灵床被大光荣环绕，宗徒高声歌唱。看，其中一人，是犹太人的大司祭，充满愤怒和狂怒，对众人说：\Quote{看，那扰乱我们和我们全族的帐幕，得到了何等的荣耀！}他上前去，想要推翻灵床，把身体扔到地上。他的双手立刻从肘部枯干，粘在了灵床上。宗徒抬起灵床时，他一部分悬着，一部分粘在灵床上；宗徒行走歌唱时，他因剧痛受折磨。云中的天使击打百姓，使他们眼瞎。\par

12. 那位首领喊叫说：\Quote{圣伯多禄，我恳求你，不要在这极大的困境中轻视我，因为我正遭受极大的折磨。你要记住，在总督府里，那看门的使女认出你\BibleRef{若 18:17}，并告诉别人辱骂你时，我曾为你说过好话。}伯多禄回答说：\Quote{我不能给你别的；但如果你全心相信主耶稣基督，就是她曾怀在胎中、并在生产后仍保持童贞的那一位，那么主的怜悯——祂以丰盛的慈爱拯救不配的人——将赐给你救恩。}\par

他回答说：\Quote{我们难道不信吗？但我们该做什么呢？人类之敌蒙蔽了我们的心，混乱遮盖了我们的脸面，免得我们承认天主的大事，尤其当我们自己曾对基督发出诅咒，喊说：\InnerQuote{祂的血归在我们和我们的儿女身上}\BibleRef{玛 27:25}。}于是伯多禄说：\Quote{看，这诅咒将伤害那对祂不忠的人；但那些转向天主的人，慈悲不会拒绝。}他说：\Quote{我相信你对我说的一切；只求你怜悯我，免得我死。}\par

13. 于是伯多禄使床榻停住，对他说：\Quote{你若是全心相信主耶稣基督，你的手就会从担架上松开。}他说了这话，他的手立刻从担架上松开，他开始站起来；但他的双臂枯干了，痛苦并没有离开他。于是伯多禄对他说：\Quote{你到遗体前，亲吻床榻，说：\InnerQuote{我相信天主，也相信天主之子耶稣基督，就是她所生的那一位，并且我相信天主的宗徒伯多禄对我所说的一切。}}他上前亲吻了床榻，立刻所有的痛苦都离开了他，他的手也痊愈了。于是他开始大大赞美天主，并从梅瑟的书引用证言来颂扬基督的赞美，以致宗徒们自己也惊奇，喜极而泣，赞美主的名。\par

14. 伯多禄对他说：\Quote{你从我们兄弟若望手中接过这棕榈枝，进城去，你会遇见许多瞎眼的人，向他们宣告天主的大事；凡信主耶稣基督的，你就把这棕榈枝放在他们的眼睛上，他们就会看见；但不信的人将仍然瞎眼。}他照做了，就遇见许多瞎眼的人，哀哭说：\Quote{我们有祸了，因为我们变得像被击打瞎眼的索多玛人一般。现在我们只有灭亡了。}但当他们听见那被治愈的首领所说的话，就相信了主耶稣基督；他把棕榈枝放在他们的眼睛上，他们就恢复了视力。其中有五人心里刚硬，死了。司祭长出城，把棕榈枝带回给宗徒们，报告了所发生的一切事。\par

15. 宗徒们抬着玛利亚，来到主曾指示他们的约沙法特谷的地方；他们将玛利亚安放在一座新坟墓里，封闭了墓穴。他们自己坐在墓门旁，正如主所命令的；忽然，主耶稣基督带着一大群天使来临，有极大光明的光环闪耀，对宗徒们说：\Quote{愿平安与你们同在！}他们回答说：\Quote{主啊，愿祢的仁慈临于我们，因为我们寄望于祢。}救主便对他们说：\Quote{在我升到我父那里之前，我曾应许你们说：你们这些在重生中跟随我的人，当人子坐在祂尊威的宝座上时，你们也要坐在十二个宝座上，审判以色列十二支派。\BibleRef{玛 19:28}因此，我凭借我父的命令，从以色列支派中拣选了她，使我居住在她内。那么，你们愿意我为她做什么？}于是伯多禄和其他宗徒说：\Quote{主，祢预先拣选了祢的这位婢女，使她成为祢无玷的居所，又拣选了我们这些仆人来服事祢。祢与父在万世之前预知一切，祢与父及圣神是同一的天主性，相等而无限的权能。因此，若在祢恩宠的权能面前可行，在祢的仆人看来，既然祢战胜了死亡，在光荣中为王，同样，求祢使祢母亲的遗体复活，高兴地把她带到祢天上的荣耀里。}\par

16. 于是救主说：\Quote{就照你们的意愿成就吧。}祂命令总领天使弥额尔将圣玛利亚的灵魂带来。看，总领天使弥额尔把石头从墓门滚开；主说：\Quote{起来，我心爱的、我最亲近的\Emph{至亲}；你未曾因与男人的交接而染上朽坏，不容许你的身体在坟墓中遭受毁坏。}玛利亚立刻从坟墓中起来，赞美主，俯伏在主脚前朝拜祂，说：\Quote{主啊，我无法充分感谢祢无限的恩惠，祢竟屈尊赐予祢的婢女这些恩惠。愿祢的名，世界的救赎者，以色列的天主，永受赞美。}\par

17. 主亲吻了她，转身回去，将她的灵魂交托给天使，让他们带入乐园。祂对宗徒们说：\Quote{到我这里来。}他们上前来，祂亲吻了他们，说：\Quote{愿平安归于你们！我如何常与你们同在，也将如此直到世界终尽。}主说了这话，立刻被云彩接去，被带回天堂，天使们也同祂一起，带着真福玛利亚进入天主的乐园。宗徒们被云彩接去，各人回到为他所指定的传道地方，讲述天主的大事，赞美我们的主耶稣基督，祂与圣父及圣神，在完美合一及同一的天主性实体中，永生永王，直到永远。阿们。\par

\SourceRecord{Translated by Alexander Walker.；Ante-Nicene Fathers；Vol. 8.；Edited by Alexander Roberts, James Donaldson, and A. Cleveland Coxe.；Buffalo, NY: Christian Literature Publishing Co.,；1886.；<http://www.newadvent.org/fathers/0832.htm>.}
